\documentclass[11pt]{article}

\usepackage[a4paper,margin=1in]{geometry}
\usepackage[T1]{fontenc}
\usepackage[utf8]{inputenc}
\usepackage{lmodern}
\usepackage{microtype}
\usepackage{amsmath,amssymb,amsfonts,mathtools}
\usepackage{amsthm}
\usepackage{enumitem}
\usepackage{booktabs}
\usepackage{graphicx}
\usepackage{float}
\usepackage{caption}
\usepackage{hyperref}

\hypersetup{
  colorlinks=true,
  linkcolor=blue,
  citecolor=blue,
  urlcolor=blue
}

\newtheorem{theorem}{Theorem}[section]
\newtheorem{proposition}[theorem]{Proposition}
\newtheorem{lemma}[theorem]{Lemma}
\newtheorem{corollary}[theorem]{Corollary}
\theoremstyle{definition}
\newtheorem{example}[theorem]{Example}
\newtheorem{definition}[theorem]{Definition}
\newtheorem{assumption}[theorem]{Assumption}
\theoremstyle{remark}
\newtheorem{remark}[theorem]{Remark}

\newcommand{\R}{\mathbb{R}}
\newcommand{\norm}[1]{\left\lVert #1\right\rVert}
\newcommand{\abs}[1]{\left\lvert #1\right\rvert}
\newcommand{\dd}{\,\mathrm{d}}

\title{A Planar Differential Inclusion with a Rotating Ellipsoidal Control Constraint}
\author{}
\date{}

\begin{document}
\maketitle

\begin{abstract}
We describe a finite-dimensional benchmark problem for controlled differential inclusions. The example consists of a nonlinear time-dependent drift in $\R^2$ and a set-valued control constraint given by a rotating, time-dependent ellipse. The admissible velocity set is represented as a Minkowski sum of the drift and the control ellipse. A smooth finite-dimensional parametrization of measurable selectors is introduced in such a way that the geometric constraint is satisfied by construction. This yields a numerically tractable optimal-control problem suitable for testing selector-based and neural parametrizations of differential inclusions.
\end{abstract}

\section{Example}
\label{sec:example-rotating-ellipse}

In this section we present a two-dimensional controlled differential inclusion with a nonautonomous drift and a rotating ellipsoidal set of admissible controls. The purpose of the example is twofold. First, it gives a transparent geometric model of a set-valued velocity field. Second, it provides a finite-dimensional numerical test problem in which admissibility of the selector is enforced by the parametrization itself, rather than through a penalty term.

Let $T>0$ and set $J=[0,T]$. In the numerical realization considered below we take
\begin{equation}
    T=3, \qquad N=120, \qquad K=24,
\end{equation}
where $N$ is the number of time-grid points and $K$ is the number of control nodes used in the selector parametrization. The prescribed initial and target states are
\begin{equation}
    x_0=(0,0)^{\top}, \qquad x_1=(1,-0.2)^{\top}.
\end{equation}
The values of the selector-regularization parameter used in the continuation procedure are
\begin{equation}
    \lambda_{\xi}\in\{10^{-3},10^{-2},10^{-1}\}.
\end{equation}

We consider trajectories $x:J\to\R^2$ satisfying the differential inclusion
\begin{equation}
    \dot{x}(t)\in F(t,x(t)), \qquad t\in J,
    \label{eq:main-inclusion}
\end{equation}
where the set-valued velocity field is of Minkowski type,
\begin{equation}
    F(t,x)=g(t,x)+E(t).
    \label{eq:minkowski-field}
\end{equation}
Equivalently, an admissible trajectory may be written as the solution of
\begin{equation}
    \dot{x}(t)=g(t,x(t))+\xi(t), \qquad \xi(t)\in E(t) \quad \text{for a.e. }t\in J.
    \label{eq:controlled-ode}
\end{equation}
Here $g$ is the drift and $\xi$ is a measurable selector of the time-dependent control set $E(t)$.

\subsection{The drift field}

The drift $g:J\times\R^2\to\R^2$ is defined by
\begin{equation}
    g(t,x)=
    \begin{pmatrix}
    \sin(x_1)+0.15\cos\left(\dfrac{2\pi t}{T}\right)\\[2mm]
    0.5\cos(x_2)-0.20\sin\left(\dfrac{2\pi t}{T}\right)
    \end{pmatrix},
    \qquad x=(x_1,x_2)^{\top}.
    \label{eq:drift}
\end{equation}
Thus the first component depends nonlinearly on $x_1$, the second component depends nonlinearly on $x_2$, and both components contain a periodic time forcing. In the geometric representation of the inclusion, the curve $t\mapsto g(t,x(t))$ forms the centerline of the admissible velocity tube.

\subsection{The rotating control ellipse}

For $t\in J$, let
\begin{equation}
    a(t)=0.3+0.4\left|\sin\left(\frac{\pi t}{T}\right)\right|,
    \qquad
    b(t)=0.15+0.25\left|\cos\left(\frac{\pi t}{T}\right)\right|,
    \label{eq:axes}
\end{equation}
and define the rotation angle
\begin{equation}
    \theta(t)=\frac{\pi t}{T}.
    \label{eq:angle}
\end{equation}
The rotation matrix is denoted by
\begin{equation}
    R_{\theta(t)}=
    \begin{pmatrix}
    \cos\theta(t)&-\sin\theta(t)\\
    \sin\theta(t)& \cos\theta(t)
    \end{pmatrix}.
    \label{eq:rotation}
\end{equation}
The admissible selector set is the ellipse
\begin{equation}
    E(t)=\left\{R_{\theta(t)}
    \begin{pmatrix}u\\w\end{pmatrix}\in\R^2:
    \frac{u^2}{a(t)^2}+\frac{w^2}{b(t)^2}\leq 1\right\}.
    \label{eq:ellipse}
\end{equation}
Since $\theta(0)=0$ and $\theta(T)=\pi$, the ellipse performs a half-turn over the interval $J$. Consequently, the family of translated sets $g(t,x(t))+E(t)$ forms a twisted tube of admissible velocities along the drift curve.

For $\xi=(\xi_1,\xi_2)^{\top}\in\R^2$, the coordinates in the local frame of the ellipse are given by
\begin{equation}
    \begin{pmatrix}u(t)\\w(t)\end{pmatrix}=R_{-\theta(t)}\xi(t),
    \label{eq:local-coordinates}
\end{equation}
that is,
\begin{equation}
    u(t)=\cos\theta(t)\xi_1(t)+\sin\theta(t)\xi_2(t),
    \qquad
    w(t)=-\sin\theta(t)\xi_1(t)+\cos\theta(t)\xi_2(t).
    \label{eq:uw-coordinates}
\end{equation}
It is convenient to introduce the ellipse-level function
\begin{equation}
    \ell(t,\xi)=
    \left(\frac{u(t)}{a(t)}\right)^2+
    \left(\frac{w(t)}{b(t)}\right)^2.
    \label{eq:level-function}
\end{equation}
Then $\xi(t)\in E(t)$ if and only if $\ell(t,\xi(t))\leq 1$.

\subsection{Selector parametrization}
\label{subsec:parametrization}

The selector is not optimized independently at all grid points. Instead, it is generated from $K$ control nodes and then interpolated to the full time grid. This reduces artificial oscillations and gives a low-dimensional representation of the admissible selector.

Let
\begin{equation}
    w_c^j\in\R^2, \qquad \rho_c^j\in\R, \qquad j=0,\ldots,K-1,
    \label{eq:control-nodes}
\end{equation}
be the nodal variables. Linear interpolation of these nodes yields functions $w(t_i)\in\R^2$ and $\rho(t_i)\in\R$ on the grid
\begin{equation}
    t_i=i\Delta t, \qquad i=0,\ldots,N-1, \qquad
    \Delta t=\frac{T}{N-1}.
    \label{eq:grid}
\end{equation}
For a small parameter $\varepsilon>0$, taken numerically as $\varepsilon=10^{-10}$, define
\begin{equation}
    d(t_i)=\frac{w(t_i)}{\sqrt{\norm{w(t_i)}^2+\varepsilon}}.
    \label{eq:direction}
\end{equation}
The radial magnitude is generated by the sigmoid function
\begin{equation}
    s(t_i)=\sigma(\rho(t_i))=\frac{1}{1+e^{-\rho(t_i)}}.
    \label{eq:sigmoid}
\end{equation}
Therefore $0<s(t_i)<1$. The local coordinates of the selector are set to be
\begin{equation}
    u(t_i)=s(t_i)a(t_i)d_1(t_i),
    \qquad
    w(t_i)=s(t_i)b(t_i)d_2(t_i),
    \label{eq:local-selector}
\end{equation}
and the selector itself is defined by
\begin{equation}
    \xi(t_i)=R_{\theta(t_i)}
    \begin{pmatrix}
    s(t_i)a(t_i)d_1(t_i)\\
    s(t_i)b(t_i)d_2(t_i)
    \end{pmatrix}.
    \label{eq:selector}
\end{equation}
Equivalently,
\begin{align}
    \xi_1(t_i)
    &=s(t_i)\left[
    \cos\theta(t_i)a(t_i)d_1(t_i)
    -\sin\theta(t_i)b(t_i)d_2(t_i)
    \right],
    \label{eq:xi-one}\\
    \xi_2(t_i)
    &=s(t_i)\left[
    \sin\theta(t_i)a(t_i)d_1(t_i)
    +\cos\theta(t_i)b(t_i)d_2(t_i)
    \right].
    \label{eq:xi-two}
\end{align}

\begin{proposition}[Admissibility by construction]
\label{prop:admissibility}
For every grid point $t_i$, the selector defined by \eqref{eq:selector} satisfies
\begin{equation}
    \xi(t_i)\in E(t_i).
\end{equation}
More precisely,
\begin{equation}
    \ell(t_i,\xi(t_i))
    =s(t_i)^2\left(d_1(t_i)^2+d_2(t_i)^2\right)
    \leq s(t_i)^2<1,
    \label{eq:level-bound}
\end{equation}
up to the harmless normalization regularization introduced by $\varepsilon$.
\end{proposition}

\begin{proof}
By construction, the local coordinates of $\xi(t_i)$ are precisely those in \eqref{eq:local-selector}. Substituting them into \eqref{eq:level-function} gives
\begin{equation}
    \ell(t_i,\xi(t_i))
    =\left(\frac{s(t_i)a(t_i)d_1(t_i)}{a(t_i)}\right)^2
    +\left(\frac{s(t_i)b(t_i)d_2(t_i)}{b(t_i)}\right)^2.
\end{equation}
This yields \eqref{eq:level-bound}. Since $0<s(t_i)<1$ and $\norm{d(t_i)}\leq 1$, the selector lies in the interior of $E(t_i)$.
\end{proof}

\begin{remark}
The constraint $\xi(t_i)\in E(t_i)$ is therefore not imposed by adding a penalty for constraint violation. It is built into the parametrization. This is useful in numerical methods for differential inclusions because the residual relation
\begin{equation}
    \dot{x}(t)-g(t,x(t))\in E(t)
\end{equation}
can be replaced by the equality
\begin{equation}
    \dot{x}(t)-g(t,x(t))=\xi(t), \qquad \xi(t)\in E(t),
\end{equation}
where the second condition is automatically satisfied.
\end{remark}

\subsection{State dynamics and time integration}

For a fixed selector $\xi$, the state equation is the ordinary differential equation
\begin{equation}
    \dot{x}(t)=g(t,x(t))+\xi(t), \qquad x(0)=x_0.
    \label{eq:state-dynamics}
\end{equation}
The numerical goal is to find an admissible selector for which
\begin{equation}
    x(T)\approx x_1.
\end{equation}
The dynamics are integrated on the grid \eqref{eq:grid} by the classical fourth-order Runge--Kutta method. Writing
\begin{equation}
    f(t,x)=g(t,x)+\xi(t),
\end{equation}
one step of the scheme is
\begin{align}
    k_1 &= f(t_k,x_k),\nonumber\\
    k_2 &= f\left(t_k+\frac{\Delta t}{2},x_k+\frac{\Delta t}{2}k_1\right),\nonumber\\
    k_3 &= f\left(t_k+\frac{\Delta t}{2},x_k+\frac{\Delta t}{2}k_2\right),\nonumber\\
    k_4 &= f(t_k+\Delta t,x_k+\Delta t k_3),\nonumber
\end{align}
and
\begin{equation}
    x_{k+1}=x_k+\frac{\Delta t}{6}\left(k_1+2k_2+2k_3+k_4\right).
    \label{eq:rk4}
\end{equation}
Between grid points, the selector is evaluated by linear interpolation.

\subsection{Finite-dimensional optimal-control problem}

The optimization variable is
\begin{equation}
    z=\left(w_c^0,\ldots,w_c^{K-1},\rho_c^0,\ldots,\rho_c^{K-1}\right)\in\R^{3K}.
    \label{eq:optimization-variable}
\end{equation}
For $K=24$ this gives $3K=72$ scalar variables. For each $z$, the selector $\xi_z(t_i)$ is reconstructed using \eqref{eq:selector}; the state $x_z(t_i)$ is obtained from \eqref{eq:rk4}; and the selected velocity is
\begin{equation}
    v_z(t_i)=g(t_i,x_z(t_i))+\xi_z(t_i).
\end{equation}
The cost functional is taken in the form
\begin{equation}
    J(z)=\lambda_TJ_T(z)+\lambda_{\xi}J_{\xi}(z)+\lambda_sJ_s(z)
    +\lambda_cJ_c(z)+\lambda_bJ_b(z)+\lambda_uJ_u(z).
    \label{eq:general-cost}
\end{equation}
The individual terms are defined as follows:
\begin{align}
    J_T(z)&=\norm{x_z(T)-x_1}^2,\label{eq:terminal-cost}\\
    J_{\xi}(z)&=\frac{1}{N}\sum_{i=0}^{N-1}\ell(t_i,\xi_z(t_i)),\label{eq:selector-energy}\\
    J_s(z)&=\frac{1}{N-1}\sum_{i=0}^{N-2}\norm{\xi_z(t_{i+1})-\xi_z(t_i)}^2,\label{eq:smoothness-cost}\\
    J_c(z)&=\frac{1}{N-2}\sum_{i=1}^{N-2}
    \norm{\xi_z(t_{i+1})-2\xi_z(t_i)+\xi_z(t_{i-1})}^2,\label{eq:curvature-cost}\\
    J_b(z)&=\frac{1}{N}\sum_{i=0}^{N-1}\ell(t_i,\xi_z(t_i))^3,\label{eq:boundary-cost}\\
    J_u(z)&=\frac{1}{K-1}\sum_{j=0}^{K-2}\norm{w_c^{j+1}-w_c^j}^2
    +\frac{1}{K-1}\sum_{j=0}^{K-2}(\rho_c^{j+1}-\rho_c^j)^2.
    \label{eq:node-smoothness}
\end{align}
The terminal weight is chosen as
\begin{equation}
    \lambda_T=2500.
\end{equation}
The term $J_T$ enforces closeness to the prescribed terminal state. The term $J_{\xi}$ discourages selectors located near the boundary of the ellipse. The terms $J_s$ and $J_c$ control first- and second-order discrete variations of the selector, while $J_b$ acts as an additional soft barrier against boundary saturation. Finally, $J_u$ regularizes the nodal variables before interpolation.

With the weights used in the implementation, the problem may be written explicitly as
\begin{align}
\min_{z\in\R^{72}}\quad
&2500\norm{x_z(T)-x_1}^2
+\lambda_{\xi}\frac{1}{N}\sum_{i=0}^{N-1}\ell(t_i,\xi_z(t_i))
\nonumber\\
&+\frac{1}{N-1}\sum_{i=0}^{N-2}\norm{\xi_z(t_{i+1})-\xi_z(t_i)}^2
\nonumber\\
&+0.2\frac{1}{N-2}\sum_{i=1}^{N-2}
\norm{\xi_z(t_{i+1})-2\xi_z(t_i)+\xi_z(t_{i-1})}^2
\nonumber\\
&+10^{-2}\frac{1}{N}\sum_{i=0}^{N-1}\ell(t_i,\xi_z(t_i))^3
\nonumber\\
&+5\cdot 10^{-2}
\left[
\frac{1}{K-1}\sum_{j=0}^{K-2}\norm{w_c^{j+1}-w_c^j}^2
+
\frac{1}{K-1}\sum_{j=0}^{K-2}(\rho_c^{j+1}-\rho_c^j)^2
\right],
\label{eq:explicit-optimization}
\end{align}
subject to
\begin{equation}
    \dot{x}_z(t)=g(t,x_z(t))+\xi_z(t), \qquad x_z(0)=x_0,
    \label{eq:optimization-dynamics}
\end{equation}
and the built-in admissibility condition
\begin{equation}
    \xi_z(t_i)\in E(t_i), \qquad i=0,\ldots,N-1.
    \label{eq:built-in-admissibility}
\end{equation}
The optimization is performed by the L-BFGS-B method, using warm starts along the sequence of regularization parameters $\lambda_{\xi}\in\{10^{-3},10^{-2},10^{-1}\}$.

\subsection{Geometric interpretation}

At each time $t_i$, the translated ellipse
\begin{equation}
    g(t_i,x_z(t_i))+E(t_i)
\end{equation}
represents the set of all admissible velocities of the inclusion. Plotting these ellipses for all grid points produces a twisted tube in the space of coordinates $(\mathrm{component\ }1,\mathrm{component\ }2,t)$. The centerline of this tube is the drift curve $g(t_i,x_z(t_i))$, while the selected velocity is
\begin{equation}
    v_z(t_i)=g(t_i,x_z(t_i))+\xi_z(t_i).
\end{equation}
Thus each vector $\xi_z(t_i)$ is the displacement from the center of the admissible velocity set to the actually selected velocity. Figure~\ref{fig:velocity-tube} shows this geometry: the drift curve forms the centerline, while the sampled velocities lie inside the translated rotating ellipses.

\begin{figure}[p]
    \centering
    \includegraphics[width=0.92\textwidth]{fig_velocity_tube.pdf}
    \caption{Admissible velocity tube associated with the inclusion $\dot{x}(t)\in g(t,x(t))+E(t)$. The blue curve represents the center $g(t,x(t))$, the orange points represent sampled selected velocities $v(t_i)=g(t_i,x(t_i))+\xi(t_i)$, and the grey segments indicate the selector displacement $\xi(t_i)=v(t_i)-g(t_i,x(t_i))$.}
    \label{fig:velocity-tube}
\end{figure}

Equivalently, in the coordinate space $(\xi_1,\xi_2,t)$ one observes the selector itself inside the time-dependent ellipse $E(t)$ centered at the origin. The value of the level function $\ell(t_i,\xi_z(t_i))$ measures the normalized radial position of the selector inside the corresponding ellipse; this is illustrated in Figure~\ref{fig:selector-cylinder}.

\begin{figure}[p]
    \centering
    \includegraphics[width=0.92\textwidth]{fig_selector_cylinder.pdf}
    \caption{Selector trajectory inside the rotating control cylinder. The surface represents the family of admissible ellipses $E(t)$, while the sampled points show $\xi(t_i)$. The color scale associated with the points represents the level $\ell(t_i,\xi(t_i))$, and the surface coloring indicates the time-dependent anisotropy ratio $a(t)/b(t)$.}
    \label{fig:selector-cylinder}
\end{figure}

The corresponding state trajectory is shown in Figure~\ref{fig:state-trajectory}. It starts at $x_0$ and approaches the prescribed terminal profile $x_1$ under the selected admissible velocities.

\begin{figure}[p]
    \centering
    \includegraphics[width=0.92\textwidth]{fig_state_trajectory.pdf}
    \caption{State trajectory $x(t)$ generated by the optimized admissible selector. The trajectory starts from $x_0=(0,0)^\top$ and is steered toward the target state $x_1=(1,-0.2)^\top$ over the interval $[0,T]$.}
    \label{fig:state-trajectory}
\end{figure}

\begin{remark}[Relation to selector-based methods]
This example can be viewed as a prototype for numerical methods for controlled differential inclusions. In the language of control theory, $\xi$ is a geometrically constrained control. In the language of differential inclusions, $\xi$ is a measurable selection of the set-valued map $t\mapsto E(t)$. In neural or physics-informed parametrizations, the same construction suggests replacing a penalty-based treatment of the inclusion constraint by an internal parametrization of admissible selectors.
\end{remark}

\begin{theorem}[Numerical feasibility of the selector formulation]
\label{thm:feasibility}
Every discrete selector generated by \eqref{eq:selector} yields a discrete trajectory satisfying the inclusion constraint at the grid level, namely
\begin{equation}
    \frac{x_{i+1}-x_i}{\Delta t}\in g(t_i,x_i)+E(t_i)
\end{equation}
up to the consistency error of the chosen time-integration scheme and the interpolation of $\xi$.
\end{theorem}

\begin{proof}
By Proposition~\ref{prop:admissibility}, $\xi(t_i)\in E(t_i)$ for every grid point. The right-hand side used in the time integrator has the form $g(t_i,x_i)+\xi(t_i)$, and therefore belongs to $g(t_i,x_i)+E(t_i)$. The stated assertion follows after accounting for the standard local truncation error of the Runge--Kutta discretization and the interpolation used to evaluate $\xi$ between grid points.
\end{proof}

In summary, the example represents the controlled inclusion
\begin{equation}
    \dot{x}(t)\in g(t,x(t))+E(t)
\end{equation}
by a smooth parametrized selector $\xi(t)\in E(t)$ and solves the finite-dimensional problem of choosing this selector so that the trajectory starting from $x_0$ reaches, as closely as possible, the prescribed target $x_1$ at time $T$.

\end{document}
